\studySubsection{practice-calc-17}{练习库：高等数学第 17 讲 多元函数积分学的预备知识（仅数学一）}

\problemAnchor{MATH1-CALC-0062}
\begin{problemBox}
\begin{problemMeta}
\textbf{题目编号：} \problemRef{MATH1-CALC-0062}\quad
\textbf{答案：} \answerRef{MATH1-CALC-0062}\quad
\textbf{来源：}Codex 原创 / K113 深度讲义配套练习（非真题）
\end{problemMeta}
\begin{enumerate}
  \item[\textbf{A}] 已知 \(P=(-1,2,3),Q=(2,-2,5)\)，求
  \(\overrightarrow{PQ}\)、\(|PQ|\) 及其同向单位向量。
  \item[\textbf{B}] 已知 \(\vec a=(\lambda,2,\lambda-1)\) 且
  \(|\vec a|=3\)，求 \(\lambda\)。
  \item[\textbf{C}] 求从 \(A=(1,1,1)\) 指向 \(B=(-1,2,3)\)
  的单位向量。
  \item[\textbf{M}] 设 \(P=(k,0,0),Q=(0,1,2)\)，已知
  \(|PQ|=3\)，且 \(\overrightarrow{PQ}\) 的同向单位向量第一分量为正，
  求 \(k\)。
\end{enumerate}
\end{problemBox}

\problemAnchor{MATH1-CALC-0063}
\begin{problemBox}
\begin{problemMeta}
\textbf{题目编号：} \problemRef{MATH1-CALC-0063}\quad
\textbf{答案：} \answerRef{MATH1-CALC-0063}\quad
\textbf{来源：}Codex 原创 / K114 深度讲义配套练习（非真题）
\end{problemMeta}
\begin{enumerate}
  \item[\textbf{A}] 设 \(\vec a=(2,-1,0),\vec b=(-1,3,2)\)，求
  \(3\vec a+2\vec b\)。
  \item[\textbf{B}] 若 \((2,\lambda,6)\parallel(1,-2,3)\)，求
  \(\lambda\)。
  \item[\textbf{C}] 求 \(m\)，使
  \(A=(0,0,0),B=(1,0,1),C=(0,2,1),D=(2,2,m)\) 共面。
  \item[\textbf{M}] 设
  \(\vec a=(1,1,0),\vec b=(0,1,1),\vec c=(1,t,2)\)。
  求 \(t\)，使 \(\vec c\) 可由 \(\vec a,\vec b\) 线性表示，
  并写出该线性表示。
\end{enumerate}
\end{problemBox}

\problemAnchor{MATH1-CALC-0064}
\begin{problemBox}
\begin{problemMeta}
\textbf{题目编号：} \problemRef{MATH1-CALC-0064}\quad
\textbf{答案：} \answerRef{MATH1-CALC-0064}\quad
\textbf{来源：}Codex 原创 / K115 深度讲义配套练习（非真题）
\end{problemMeta}
\begin{enumerate}
  \item[\textbf{A}] 求 \((1,-1,0)\) 与 \((0,1,-1)\) 的夹角。
  \item[\textbf{B}] 求 \(\vec a=(1,3,2)\) 在
  \(\vec b=(2,0,1)\) 方向的标量投影和向量投影。
  \item[\textbf{C}] 求 \(\lambda\)，使
  \((\lambda,1,2)\perp(1,-2,1)\)。
  \item[\textbf{M}] 求实数 \(t\)，使
  \(|(2,1,2)-t(1,1,0)|\) 最小，并求该最小值。
\end{enumerate}
\end{problemBox}

\problemAnchor{MATH1-CALC-0065}
\begin{problemBox}
\begin{problemMeta}
\textbf{题目编号：} \problemRef{MATH1-CALC-0065}\quad
\textbf{答案：} \answerRef{MATH1-CALC-0065}\quad
\textbf{来源：}Codex 原创 / K116 深度讲义配套练习（非真题）
\end{problemMeta}
\begin{enumerate}
  \item[\textbf{A}] 计算
  \((1,-1,2)\times(2,0,1)\)。
  \item[\textbf{B}] 求三角形
  \(A=(1,0,0),B=(0,2,0),C=(0,0,3)\) 的面积。
  \item[\textbf{C}] 求所有同时垂直于
  \((1,0,1)\)、\((0,2,1)\) 的单位向量。
  \item[\textbf{M}] 已知 \((1,t,1)\) 与 \((1,1,-1)\)
  张成的平行四边形面积为 \(\sqrt6\)，求 \(t\)。
\end{enumerate}
\end{problemBox}

\problemAnchor{MATH1-CALC-0066}
\begin{problemBox}
\begin{problemMeta}
\textbf{题目编号：} \problemRef{MATH1-CALC-0066}\quad
\textbf{答案：} \answerRef{MATH1-CALC-0066}\quad
\textbf{来源：}Codex 原创 / K117 深度讲义配套练习（非真题）
\end{problemMeta}
\begin{enumerate}
  \item[\textbf{A}] 计算
  \([(1,1,0),(0,1,2),(2,0,1)]\)。
  \item[\textbf{B}] 求四面体
  \(O=(0,0,0),A=(2,0,0),B=(0,3,0),C=(1,1,2)\) 的体积。
  \item[\textbf{C}] 求 \(t\)，使
  \((1,0,1),(0,2,1),(1,t,2)\) 共面。
  \item[\textbf{M}] 求四面体
  \(A=(1,0,0),B=(0,1,0),C=(0,0,1),D=(1,1,1)\) 的体积，
  并说明交换任意两个顶点为何不改变几何体积。
\end{enumerate}
\end{problemBox}

\problemAnchor{MATH1-CALC-0067}
\begin{problemBox}
\begin{problemMeta}
\textbf{题目编号：} \problemRef{MATH1-CALC-0067}\quad
\textbf{答案：} \answerRef{MATH1-CALC-0067}\quad
\textbf{来源：}Codex 原创 / K118 深度讲义配套练习（非真题）
\end{problemMeta}
\begin{enumerate}
  \item[\textbf{A}] 求方向数 \((1,-2,2)\) 对应的方向余弦。
  \item[\textbf{B}] 已知 \(\cos\alpha=\cos\beta=1/2\)，且
  \(\gamma\) 为钝角，求 \(\cos\gamma\)。
  \item[\textbf{C}] 某有向直线的方向余弦之比为 \(-1:2:2\)，
  求其方向余弦。
  \item[\textbf{M}] 某有向直线满足 \(\alpha=\beta\)、
  \(\gamma=60^\circ\)，且 \(\cos\alpha,\cos\beta>0\)。
  求三个方向余弦。
\end{enumerate}
\end{problemBox}

\problemAnchor{MATH1-CALC-0068}
\begin{problemBox}
\begin{problemMeta}
\textbf{题目编号：} \problemRef{MATH1-CALC-0068}\quad
\textbf{答案：} \answerRef{MATH1-CALC-0068}\quad
\textbf{来源：}Codex 原创 / K119 深度讲义配套练习（非真题）
\end{problemMeta}
\begin{enumerate}
  \item[\textbf{A}] 求过 \((2,-1,1)\)、法向量为 \((1,2,-2)\)
  的平面方程。
  \item[\textbf{B}] 求过
  \(A=(0,0,1),B=(1,0,0),C=(0,2,0)\) 的平面。
  \item[\textbf{C}] 求包含直线
  \(\vec r=(1,0,0)+t(1,1,0)\) 且过点 \(Q=(0,0,1)\) 的平面。
  \item[\textbf{M}] 求过原点且同时垂直于平面
  \(x+y=0\)、\(y+z=0\) 的平面。
\end{enumerate}
\end{problemBox}

\problemAnchor{MATH1-CALC-0069}
\begin{problemBox}
\begin{problemMeta}
\textbf{题目编号：} \problemRef{MATH1-CALC-0069}\quad
\textbf{答案：} \answerRef{MATH1-CALC-0069}\quad
\textbf{来源：}Codex 原创 / K120 深度讲义配套练习（非真题）
\end{problemMeta}
\begin{enumerate}
  \item[\textbf{A}] 写出过 \((-1,2,0)\)、方向为 \((2,-1,3)\)
  的直线参数式与对称式。
  \item[\textbf{B}] 写出过 \((1,0,-2)\)、方向为 \((0,1,2)\)
  的直线，且不得出现零分母。
  \item[\textbf{C}] 求平面 \(x+y=1\) 与 \(y+z=2\) 的交线。
  \item[\textbf{M}] 求过 \(P=(1,2,0)\)，且平行于平面
  \(x+y+z=0\) 与 \(x-y+2z=0\) 交线的直线。
\end{enumerate}
\end{problemBox}

\problemAnchor{MATH1-CALC-0070}
\begin{problemBox}
\begin{problemMeta}
\textbf{题目编号：} \problemRef{MATH1-CALC-0070}\quad
\textbf{答案：} \answerRef{MATH1-CALC-0070}\quad
\textbf{来源：}Codex 原创 / K121 深度讲义配套练习（非真题）
\end{problemMeta}
\begin{enumerate}
  \item[\textbf{A}] 判断
  \(L_1=t(1,2,0)\) 与
  \(L_2=(0,0,1)+s(2,4,0)\) 的关系。
  \item[\textbf{B}] 判断
  \(L_1=t(1,0,1)\) 与
  \(L_2=(1,0,1)+s(0,1,1)\) 的关系，并求交点（若存在）。
  \item[\textbf{C}] 判断
  \(L_1=t(1,0,0)\) 与
  \(L_2=(0,1,2)+s(0,1,0)\) 的关系。
  \item[\textbf{M}] 分别判断平面
  \[
  \Pi_1:x+2y-z=1,\quad
  \Pi_2:2x+4y-2z=3,\quad
  \Pi_3:2x-y=0
  \]
  中 \(\Pi_1,\Pi_2\) 及 \(\Pi_1,\Pi_3\) 的关系。
\end{enumerate}
\end{problemBox}

\problemAnchor{MATH1-CALC-0071}
\begin{problemBox}
\begin{problemMeta}
\textbf{题目编号：} \problemRef{MATH1-CALC-0071}\quad
\textbf{答案：} \answerRef{MATH1-CALC-0071}\quad
\textbf{来源：}Codex 原创 / K122 深度讲义配套练习（非真题）
\end{problemMeta}
\begin{enumerate}
  \item[\textbf{A}] 判断直线
  \(\vec r=(1,1,0)+t(1,-1,0)\) 与平面 \(x+y+z=2\) 的关系。
  \item[\textbf{B}] 判断直线
  \(\vec r=t(1,-1,0)\) 与平面 \(x+y+z=1\) 的关系。
  \item[\textbf{C}] 求方向为 \((1,2,-1)\) 的直线与平面
  \(2x-y+2z=0\) 的夹角。
  \item[\textbf{M}] 方向为 \((1,k,1)\) 的直线与平面 \(z=0\)
  的夹角为 \(30^\circ\)，求 \(k\)。
\end{enumerate}
\end{problemBox}

\problemAnchor{MATH1-CALC-0072}
\begin{problemBox}
\begin{problemMeta}
\textbf{题目编号：} \problemRef{MATH1-CALC-0072}\quad
\textbf{答案：} \answerRef{MATH1-CALC-0072}\quad
\textbf{来源：}Codex 原创 / K123 深度讲义配套练习（非真题）
\end{problemMeta}
\begin{enumerate}
  \item[\textbf{A}] 求点 \((2,1,-1)\) 到平面
  \(x-2y+2z+1=0\) 的距离。
  \item[\textbf{B}] 求点 \((1,1,2)\) 到 \(z\) 轴的距离。
  \item[\textbf{C}] 求点 \(P=(1,0,0)\) 在平面
  \(x+y+z=0\) 上的垂足及距离。
  \item[\textbf{M}] 点 \(P_t=(t,0,0)\) 到平面
  \(x+y+z=3\) 的距离等于它到 \(z\) 轴的距离，求 \(t\)。
\end{enumerate}
\end{problemBox}

\problemAnchor{MATH1-CALC-0073}
\begin{problemBox}
\begin{problemMeta}
\textbf{题目编号：} \problemRef{MATH1-CALC-0073}\quad
\textbf{答案：} \answerRef{MATH1-CALC-0073}\quad
\textbf{来源：}Codex 原创 / K124 深度讲义配套练习（非真题）
\end{problemMeta}
\begin{enumerate}
  \item[\textbf{A}] 求平面 \(x-2y+2z=1\) 与
  \(2x-4y+4z=-3\) 的距离。
  \item[\textbf{B}] 求平行直线
  \(L_1=t(1,-1,1)\) 与
  \(L_2=(1,2,0)+s(2,-2,2)\) 的距离。
  \item[\textbf{C}] 求异面直线
  \(L_1=t(1,0,0)\) 与
  \(L_2=(0,2,3)+s(0,1,0)\) 的距离。
  \item[\textbf{M}] 设
  \(L_1=t(1,0,0)\)，
  \(L_2=(0,a,a)+s(0,1,0)\)。
  求两直线距离为 \(2\) 时的 \(a\)，并判断此时两直线关系。
\end{enumerate}
\end{problemBox}

\problemAnchor{MATH1-CALC-0074}
\begin{problemBox}
\begin{problemMeta}
\textbf{题目编号：} \problemRef{MATH1-CALC-0074}\quad
\textbf{答案：} \answerRef{MATH1-CALC-0074}\quad
\textbf{来源：}Codex 原创 / K125 深度讲义配套练习（非真题）
\end{problemMeta}
\begin{enumerate}
  \item[\textbf{A}] 区分 \(x^2+y^2+z^2=9\) 与
  \(x^2+y^2+z^2\le9\) 所表示的几何对象。
  \item[\textbf{B}] 用截痕法说明曲面 \(z=2x^2+y^2\) 的顶点、
  开口方向和水平截面。
  \item[\textbf{C}] 说明
  \(x^2+2y^2+z^2=4\) 与
  \(\{x^2+2y^2+z^2=4,\ z=x-y\}\) 在三维中分别表示什么。
  \item[\textbf{M}] 识别曲面 \(z=4-x^2-y^2\)，给出
  \(z=c\) 截痕的存在范围，并写出它在 \(xy\) 平面的投影。
\end{enumerate}
\end{problemBox}

\problemAnchor{MATH1-CALC-0075}
\begin{problemBox}
\begin{problemMeta}
\textbf{题目编号：} \problemRef{MATH1-CALC-0075}\quad
\textbf{答案：} \answerRef{MATH1-CALC-0075}\quad
\textbf{来源：}Codex 原创 / K126 深度讲义配套练习（非真题）
\end{problemMeta}
\begin{enumerate}
  \item[\textbf{A}] 求 \(xy\) 平面内线段
  \(y=2,\ -1\le x\le1\) 绕 \(x\) 轴旋转所得曲面。
  \item[\textbf{B}] 求上半圆 \((x-1)^2+y^2=4,\ y\ge0\)
  绕 \(x\) 轴旋转所得曲面。
  \item[\textbf{C}] 求曲线 \(x=1+y^2\) 绕 \(y\) 轴旋转所得曲面。
  \item[\textbf{M}] 求曲线 \(y=\sqrt{x-1},\ x\ge1\)
  绕 \(x\) 轴旋转所得曲面，并说明平方是否引入额外分支。
\end{enumerate}
\end{problemBox}

\problemAnchor{MATH1-CALC-0076}
\begin{problemBox}
\begin{problemMeta}
\textbf{题目编号：} \problemRef{MATH1-CALC-0076}\quad
\textbf{答案：} \answerRef{MATH1-CALC-0076}\quad
\textbf{来源：}Codex 原创 / K127 深度讲义配套练习（非真题）
\end{problemMeta}
\begin{enumerate}
  \item[\textbf{A}] 识别
  \(\frac{x^2}{4}+y^2=1\) 的准线、母线方向和半轴。
  \item[\textbf{B}] 识别 \(xz=1\) 的准线与母线方向。
  \item[\textbf{C}] 识别 \(y=z^2+1\) 的准线、母线方向、顶线和开口方向。
  \item[\textbf{M}] 对集合
  \[
  D=\{(x,y,z):x^2+z^2\le4,\ -1\le y\le2\},
  \]
  指出其侧面、两个底面及内部投影，区分等式曲面与不等式实体。
\end{enumerate}
\end{problemBox}

\problemAnchor{MATH1-CALC-0077}
\begin{problemBox}
\begin{problemMeta}
\textbf{题目编号：} \problemRef{MATH1-CALC-0077}\quad
\textbf{答案：} \answerRef{MATH1-CALC-0077}\quad
\textbf{来源：}Codex 原创 / K128 深度讲义配套练习（非真题）
\end{problemMeta}
\begin{enumerate}
  \item[\textbf{A}] 识别并写出中心、半轴：
  \[
  x^2+4y^2+z^2-2x+4y=10.
  \]
  \item[\textbf{B}] 识别
  \[
  (x-1)^2+\frac{(y+1)^2}{4}-\frac{z^2}{9}=1
  \]
  并说明中心与轴向。
  \item[\textbf{C}] 识别
  \[
  \frac{z^2}{9}-\frac{x^2}{4}-y^2=1
  \]
  并给出横截面存在范围。
  \item[\textbf{M}] 识别
  \((x-1)^2+4y^2=(z+2)^2\)，并给出顶点、轴向及
  \(z=c\) 截痕的存在范围与类型。
\end{enumerate}
\end{problemBox}

\problemAnchor{MATH1-CALC-0078}
\begin{problemBox}
\begin{problemMeta}
\textbf{题目编号：} \problemRef{MATH1-CALC-0078}\quad
\textbf{答案：} \answerRef{MATH1-CALC-0078}\quad
\textbf{来源：}Codex 原创 / K129 深度讲义配套练习（非真题）
\end{problemMeta}
\begin{enumerate}
  \item[\textbf{A}] 将
  \(x=t^2,\ y=t^3,\ z=t,\ -1\le t\le2\)
  化为一般方程并保留范围。
  \item[\textbf{B}] 参数化曲线
  \[
  \frac{x^2}{4}+y^2=1,\qquad z=xy,
  \]
  使其恰描一周。
  \item[\textbf{C}] 对参数曲线
  \(\vec r(t)=(t^2-1,2t,t^3)\)，求 \(t=1\) 处一点和切向量。
  \item[\textbf{M}] 参数化两曲面的交线
  \[
  x^2+y^2+z^2=2,\qquad z=x+y,
  \]
  并说明参数区间如何避免重复端点。
\end{enumerate}
\end{problemBox}

\problemAnchor{MATH1-CALC-0079}
\begin{problemBox}
\begin{problemMeta}
\textbf{题目编号：} \problemRef{MATH1-CALC-0079}\quad
\textbf{答案：} \answerRef{MATH1-CALC-0079}\quad
\textbf{来源：}Codex 原创 / K130 深度讲义配套练习（非真题）
\end{problemMeta}
\begin{enumerate}
  \item[\textbf{A}] 求螺旋线
  \(\vec r(t)=(3\cos t,3\sin t,t)\) 在 \(xy\) 平面的投影曲线。
  \item[\textbf{B}] 求曲线
  \[
  x^2+y^2=z,\qquad z=2x
  \]
  在 \(xy\) 平面的投影曲线。
  \item[\textbf{C}] 求立体
  \[
  D=\{(x,y,z):x^2+y^2\le z\le2x\}
  \]
  在 \(xy\) 平面的投影域。
  \item[\textbf{M}] 曲线
  \[
  x=z^2+1,\qquad y=2z,\qquad -1\le z\le1
  \]
  在 \(xy\) 平面的投影是什么？须给出完整范围。
\end{enumerate}
\end{problemBox}

\problemAnchor{MATH1-CALC-0080}
\begin{problemBox}
\begin{problemMeta}
\textbf{题目编号：} \problemRef{MATH1-CALC-0080}\quad
\textbf{答案：} \answerRef{MATH1-CALC-0080}\quad
\textbf{来源：}Codex 原创 / K131 深度讲义配套练习（非真题）
\end{problemMeta}
\begin{enumerate}
  \item[\textbf{A}] 判断下列任务应留在本讲还是路由到线性代数：
  （i）配方识别 \(x^2-2x+y^2+z^2=3\)；
  （ii）自行求正交变换消去一般式中的 \(xy,xz,yz\) 项。
  \item[\textbf{B}] 不使用矩阵分类，识别
  \[
  4x^2+y^2+z^2-8x=1.
  \]
  \item[\textbf{C}] 已给
  \(u=(x+y)/\sqrt2,\ v=(x-y)/\sqrt2\)，
  化简 \(x^2-y^2\)，并判断新表达式是否出现 \(uv\) 交叉项。
  \item[\textbf{M}] 评价命题：
  “一般二次曲面主轴变换不是高数常规要求，所以数学一无需复习二次型。”
  给出正确的复习分流。
\end{enumerate}
\end{problemBox}
